\documentclass[12pt]{article}

\usepackage[utf8]{inputenc}
\usepackage{amsmath, amssymb, amsthm}
\usepackage{geometry}
\geometry{a4paper, margin=2.5cm}

\title{Teorema Banach--Steinhaus}
\author{}
\date{}

\theoremstyle{definition}
\newtheorem{definition}{Definiție}

\theoremstyle{plain}
\newtheorem{theorem}{Teoremă}
\newtheorem{proposition}{Propoziție}

\begin{document}

\maketitle

\begin{definition}
Fie $X$ un spațiu Banach și $Y$ un spațiu normat. O familie $\mathcal{F} \subset \mathcal{L}(X,Y)$
de operatori liniari și continui se numește \emph{punctat mărginită} dacă pentru fiecare $x \in X$ există o constantă $M_x > 0$ astfel încât


\[
\sup_{T \in \mathcal{F}} \|T x\| \le M_x.
\]


\end{definition}

\begin{theorem}[Banach--Steinhaus, Principiul Uniformei Limitări]
Fie $X$ un spațiu Banach și $Y$ un spațiu normat. Dacă $\mathcal{F} \subset \mathcal{L}(X,Y)$ este o familie de operatori liniari continui, punctat mărginită, atunci:


\[
\sup_{T \in \mathcal{F}} \|T\| < \infty.
\]


Cu alte cuvinte, punctata mărginitate implică mărginitate uniformă a normelor operatorilor.
\end{theorem}

\begin{proposition}
În ipotezele teoremei Banach--Steinhaus, următoarele afirmații sunt echivalente:
\begin{enumerate}
    \item[(i)] Familia $\mathcal{F}$ este punctat mărginită.
    \item[(ii)] Pentru orice mulțime mărginită $B \subset X$, mulțimea
    

\[
    \mathcal{F}(B) = \{ T(x) : T \in \mathcal{F},\, x \in B \}
    \]


    este mărginită în $Y$.
\end{enumerate}
\end{proposition}

\end{document}
